% ============================================================
% s360k.tex - 21 U.S.C. § 360k (FD&C Act § 521): State and local requirements;
% device preemption. Conforming amendment: add a new subsection (c) savings
% clause preserving State human-in-the-loop and audit requirements of general
% applicability to the extent they are not different from, or in addition to, a
% federal device requirement, marking the boundary without expanding preemption.
% NOTE: this is § 360k (State preemption), distinct from § 360(k) (510(k)
% premarket notification), reproduced in s360.tex.
% ============================================================
\uscsection{§~360k.}{State and local requirements respecting devices}

\uscprov{1}{\textbf{(a)}\textbf{ General rule} Except as provided in subsection
(b), no State or political subdivision of a State may establish or continue in
effect with respect to a device intended for human use any requirement-}
\uscprov{2}{\textbf{(1)} which is different from, or in addition to, any
requirement applicable under this chapter to the device, and}
\uscprov{2}{\textbf{(2)} which relates to the safety or effectiveness of the
device or to any other matter included in a requirement applicable to the device
under this chapter.}
\uscprov{1}{\textbf{(b)}\textbf{ Exempt requirements} Upon application of a State
or a political subdivision thereof, the Secretary may, by regulation promulgated
after notice and opportunity for an oral hearing, exempt from subsection (a),
under such conditions as may be prescribed in such regulation, a requirement of
such State or political subdivision applicable to a device intended for human use
if-}
\uscprov{2}{\textbf{(1)} the requirement is more stringent than a requirement
under this chapter which would be applicable to the device if an exemption were
not in effect under this subsection; or}
\uscprov{2}{\textbf{(2)} the requirement-}
\uscprov{3}{\textbf{(A)} is required by compelling local conditions, and}
\uscprov{3}{\textbf{(B)} compliance with the requirement would not cause the
device to be in violation of any applicable requirement under this chapter.}
\uscprov{1}{\cprinsert{\textbf{(c)}\textbf{ Rule of construction for Physical AI
oversight} Nothing in this section preempts a State requirement of general
applicability that a licensed human retain the authority to monitor, intervene,
override, or terminate an artificial intelligence or robotic system, or that an
audit record be maintained, to the extent that the requirement is not different
from, or in addition to, a requirement applicable under this chapter to a
device, including the verification requirements of section 515D.}}
